\chapter{Inferência de tipos}

\section{Introdução}

Os capítulos anteriores apresentaram a \textbf{verificação de tipos}:
dado um programa com anotações explícitas de tipo, decidir se ele é
bem tipado. A verificação percorre a árvore sintática e, para cada
nó, confirma que os tipos declarados pelo programador são
consistentes com as regras do sistema de tipos.

A \textbf{inferência de tipos} resolve um problema mais ambicioso:
dado um programa \emph{sem} anotações, encontrar automaticamente os
tipos mais gerais que tornam o programa bem tipado, ou reportar
que nenhum tal conjunto de tipos existe. O programador não precisa
anotar nada; o compilador deduz tudo.

A diferença fundamental entre os dois problemas é ilustrada pela
regra (T-Abs) do $\lambda$-cálculo simplesmente tipado:

\[ \frac{\Gamma,\, x : T_1 \vdash t : T_2} {\Gamma \vdash \lambda x {:} T_1.\;
t : T_1 \to T_2} \quad\text{(T-Abs)} \]

Na verificação, \(T_1\) é fornecido pelo programador na anotação
\(\lambda x \{:\} T\_1\). Na inferência,
\(T_1\) é desconhecido e deve ser descoberto a partir de como \(x\) é
usado no corpo \(t\).

Este capítulo estuda a inferência de tipos para a linguagem
\textbf{MiniML}, um núcleo funcional com inteiros, booleanos,
lambdas e definições polimórficas. O algoritmo é baseado em
\textbf{geração e resolução de restrições}, seguindo a abordagem
HM(X) de Pottier e Rémy. Em vez de computar os tipos diretamente, o
gerador de restrições produz um sistema de equações entre tipos; o
resolvedor então encontra a substituição, um mapeamento de
variáveis de tipo para tipos concretos, que satisfaz todas as equações.

\section{A Linguagem MiniML}

\subsection{Sintaxe}

MiniML é um fragmento de ML com cinco construções:

\[
  \begin{array}{rcll} e & ::= & x & \text{variável} \\ & | & \ell &
    \text{literal (inteiro ou booleano)} \\ & | & e_1\; e_2 & \text{aplicação}
    \\ & | & \lambda x.\; e & \text{abstração (sem anotação)} \\ & | &
    \mathbf{let}\; x = e_1\; \mathbf{in}\; e_2 & \text{def. polimórfica}
\end{array} \]

Note que as abstrações \textbf{não} carregam anotações de tipo; a
responsabilidade de descobri-las é do inferidor.

\subsection{Tipos e Esquemas}\label{tipos-e-esquemas}

Os \textbf{tipos monomórficos} (monotipos) são:

\[
  \begin{array}{rcll} T & ::= & \mathbf{Int} & \text{inteiros} \\ & | &
    \mathbf{Bool} & \text{booleanos} \\ & | & T \to T & \text{funções} \\ & | &
    \alpha & \text{variáveis de tipo}
\end{array} \]

Os \textbf{esquemas de tipo} (tipos polimórficos) generalizam
tipos monomórficos com quantificação universal:

\[ \sigma ::= T \;\mid\; \forall \vec{\alpha}.\; T \]

Escrevemos \(\forall \alpha_1 \ldots \alpha_n.\; T\) (ou simplesmente
\(\forall \vec{\alpha}.\; T\)) para o esquema que generaliza \(T\)
nas variáveis \(\alpha_1, \ldots, \alpha_n\). Um esquema com nenhuma
variável quantificada, \(\forall \varepsilon.\; T\), é equivalente ao
tipo monomórfico \(T\).

Uma \textbf{substituição} é uma função finita de variáveis de tipo
para tipos, \(S : \mathit{TVar} \rightharpoonup
\mathit{Type}\). A sua aplicação é estendida a tipos,
esquemas e contextos:

\[
  \begin{array}{lcl}
    S(\alpha) & = & \left\{
      \begin{array}{ll}
        T & \mathrm{se}\, \alpha\mapsto T \in S \\
        \alpha & \mathrm{caso contrário}\\
      \end{array}
      \right.\\
      S(T_1 \to T_2) & = & S(T_1) \to S(T_2)\\
      S(\forall \vec{\alpha}.\; T) & = & \forall \vec{\alpha}.\; S'(T) \quad
      \mathrm{onde }\: S' = S \setminus \{\alpha_1, \ldots, \alpha_n\}
    \end{array}
  \]

  A \textbf{composição} \(S_1 \circ S_2\) é definida por
  \((S\_1 \circ S\_2)(\alpha) = S\_1(S\_2(\alpha))\).

  O contexto de tipos \(\Gamma : \mathit{Var} \rightharpoonup
  \mathit{Scheme}\) associa cada variável de programa ao seu esquema.
  A \textbf{instanciação} de um esquema \(\forall
  \vec{\alpha}.\; T\) substitui as variáveis
  quantificadas por novos tipos: \([\vec{\alpha} \mapsto
  \vec{\beta}]\, T\), onde \(\vec{\beta}\) são variáveis novas. A
  \textbf{generalização} de um tipo \(T\) em um contexto \(\Gamma\)
  produz o esquema \(\forall \vec{\alpha}.\; T\), onde \(\vec{\alpha} =
  \mathrm{fv}(T) \setminus \mathrm{fv}(\Gamma)\) são as variáveis de
  tipo livres em \(T\) mas não em \(\Gamma\).

  \section{Restrições}

  \subsection{Sintaxe}

  A geração de restrições transforma o problema de inferência em um
  sistema de equações entre tipos. O conjunto de restrições
  \(\mathcal{C}\) é definido pela seguinte gramática:

  \[
    \begin{array}{rcll} C & ::= & \top & \text{restrição trivial} \\ & | & T_1
      = T_2 & \text{equação de tipos} \\ & | & C_1 \wedge C_2 & \text{conjunção}
      \\ & | & \exists \alpha.\; C & \text{existencial (variável
      local)} \\ & | &
      \mathbf{def}\; x : T\; \mathbf{in}\; C & \text{definição
      monomórfica} \\ & | &
      \mathbf{let}\; x : T\; [C_1]\; \mathbf{in}\; C_2 & \text{definição
      polimórfica}
      \\ & | & \mathbf{inst}(x,\, T) & \text{instanciação de variável} \\ & | &
      \mathbf{instScheme}(\sigma,\, T) & \text{verificação de anotação de tipo}
    \end{array}
  \]

  Cada construção tem um papel preciso no processo de inferência:

  \begin{itemize}
    \item
      \(\top\) representa a ausência de restrição, qualquer
      substituição a satisfaz.
    \item
      \(T_1 = T_2\) é uma equação de unificação: a solução deve associar
      as variáveis de tipo de forma que \(T_1\) e \(T_2\) se tornem iguais.
    \item
      \(C_1 \wedge C_2\) exige que ambas as restrições sejam
      satisfeitas simultaneamente.
    \item
      \(\exists \alpha.\; C\) introduz \(\alpha\) como uma variável de
      tipo local, presente em \(C\) mas não visível no exterior.
      Usado para nomear tipos intermediários.
    \item
      \(\mathbf{def}\; x : T\; \mathbf{in}\; C\) torna \(x\) disponível
      em \(C\) com o tipo \(T\) (não polimórfico). Usado para
      parâmetros de lambda.
    \item
      \(\mathbf{let}\; x : T\; [C_1]\; \mathbf{in}\; C_2\) primeiro
      resolve \(C_1\) para inferir o tipo de \(x\), generaliza \(T\) em
      um esquema, e então disponibiliza \(x\) com esse esquema
      polimórfico em \(C_2\).
    \item
      \(\mathbf{inst}(x,\, T)\) exige que o tipo de \(x\) (um esquema)
      seja instanciado e unificado com \(T\).
    \item
      \(\mathbf{instScheme}(\sigma,\, T)\) verifica que uma anotação de
      tipo explícita \(\sigma\) é compatível com o tipo inferido \(T\).
  \end{itemize}

  \subsection{Semântica}

  A satisfazibilidade de uma restrição é definida em relação a um
  \textbf{par} composto por uma substituição \(S\) e um contexto
  \(\Gamma\). Escrevemos \((S, \Gamma) \models
  C\) para expressar que \(C\) é satisfeita.

  \[ (S, \Gamma) \models \top \quad \text{sempre} \]

  \[ (S, \Gamma) \models T_1 = T_2 \iff S(T_1) = S(T_2) \]

  \[ (S, \Gamma) \models C_1 \wedge C_2 \iff (S, \Gamma) \models C_1 \text{ e }
  (S, \Gamma) \models C_2 \]

  \[ (S, \Gamma) \models \exists \alpha.\; C \iff \exists T.\; (S[\alpha
  \mapsto T], \Gamma) \models C \]

  \[ (S, \Gamma) \models \mathbf{def}\; x : T\; \mathbf{in}\; C \iff (S,
  \Gamma[x \mapsto \overline{\forall \varepsilon}.\; T]) \models C \]

  \[ (S, \Gamma) \models \mathbf{let}\; x : T\; [C_1]\; \mathbf{in}\; C_2
    \iff (S, \Gamma) \models C_1 \text{ e } (S, \Gamma[x \mapsto
  \mathrm{gen}(\Gamma, S(T))]) \models C_2 \]

  \[ (S, \Gamma) \models \mathbf{inst}(x,\, T) \iff \Gamma(x) = \sigma \text{ e
  } S(T) \text{ é instância de } S(\sigma) \]

  onde \(\mathrm{gen}(\Gamma, T) = \forall \vec{\alpha}.\; T\) com
  \(\vec{\alpha} = \mathrm{fv}(T) \setminus \mathrm{fv}(\Gamma)\). A
  condição \(\overline{\forall \varepsilon}.\; T\) denota o esquema sem
  variáveis quantificadas, ou seja, o monotipo \(T\).

  Uma restrição \(C\) é \textbf{satisfazível} se existe algum \((S,
  \Gamma)\) tal que \((S, \Gamma) \models C\). Uma substituição \(S\) é
  uma \textbf{solução} de \(C\) (no contexto inicial vazio) quando
  \((S, \emptyset) \models C\).

  \subsection{Geração de Restrições}

  O gerador de restrições é uma função \(\mathcal{G}(e,\, \alpha)\) que
  recebe uma expressão \(e\) e um tipo \emph{alvo} \(\alpha\) (uma
  variável de tipo fresca que representará o tipo de \(e\)) e produz
  uma restrição \(C\) tal que qualquer solução \(S\) de \(C\) satisfaz
  \(S(\alpha) = T\), onde \(T\) é o tipo de \(e\) segundo o sistema de
  tipos de Hindley-Milner.

  As regras de geração são:

  \[ \mathcal{G}(x,\; \alpha) \;=\; \mathbf{inst}(x,\, \alpha)
  \quad\text{(G-Var)} \]

  \[ \mathcal{G}(\ell,\; \alpha) \;=\; \alpha = \mathrm{tylit}(\ell)
  \quad\text{(G-Lit)} \]

  onde \(\mathrm{tylit}(n) = \mathbf{Int}\) e
  \(\mathrm{tylit}(b) = \mathbf{Bool}\).

  \[ \mathcal{G}(e_1\; e_2,\; \alpha) \;=\; \exists
    \beta.\;\bigl(\mathcal{G}(e_1,\; \beta \to \alpha) \wedge \mathcal{G}(e_2,\;
  \beta)\bigr) \quad\text{(G-App)} \]

  \[ \mathcal{G}(\lambda x.\; e,\; \alpha) \;=\; \exists \beta_1.\; \exists
    \beta_2.\; \bigl(\mathbf{def}\; x : \beta_1\; \mathbf{in}\;
    \mathcal{G}(e,\, \beta_2)\bigr) \;\wedge\; (\alpha = \beta_1 \to \beta_2)
  \quad\text{(G-Lam)} \]

  \[ \mathcal{G}(\mathbf{let}\; x = e_1\; \mathbf{in}\; e_2,\; \alpha)
    \;=\; \exists \beta.\; \mathbf{let}\; x : \beta\; [\mathcal{G}(e_1,\,
  \beta)]\; \mathbf{in}\; \mathcal{G}(e_2,\, \alpha) \quad\text{(G-Let)} \]

  Quando a expressão carrega uma anotação de esquema explícita
  \(\mathbf{let}\; x \{:\} \sigma = e\_1;
  \mathbf{in}\; e\_2\), a restrição do corpo
  da ligação recebe a conjunção adicional
  \(\mathbf{instScheme}(\sigma,\, \beta)\), verificando que o tipo
  inferido é compatível com a anotação:

  \[ \mathcal{G}(\mathbf{let}\; x {:} \sigma = e_1\; \mathbf{in}\; e_2,\;
    \alpha) \;=\; \exists \beta.\; \mathbf{let}\; x : \beta\;
    [\mathcal{G}(e_1,\, \beta) \wedge \mathbf{instScheme}(\sigma,\, \beta)]\;
  \mathbf{in}\; \mathcal{G}(e_2,\, \alpha) \quad\text{(G-LetAnn)} \]

  Para iniciar a inferência de uma expressão fechada \(e\), cria-se uma
  nova variável \(\alpha_0\) e resolve-se a restrição:

  \[ \exists \alpha_0.\; \mathcal{G}(e,\; \alpha_0) \]

  A variável \(\alpha_0\) captura o tipo principal de \(e\) após a resolução.

  \textbf{Exemplo.} Considere o termo \(\lambda f.\; f\;
  \mathbf{true}\) com alvo \(\alpha_0\). Aplicando as regras:

  \[ \mathcal{G}(\lambda f.\; f\; \mathbf{true},\; \alpha_0) = \exists
    \beta_1.\; \exists \beta_2.\; \bigl(\mathbf{def}\; f : \beta_1\;
      \mathbf{in}\; \exists \beta_3.\; \bigl(\mathcal{G}(f,\, \beta_3
        \to \beta_2)
    \wedge \mathcal{G}(\mathbf{true},\, \beta_3)\bigr) \bigr) \;\wedge\;
  (\alpha_0 = \beta_1 \to \beta_2) \]

  Expandindo \(\mathcal{G}(f, \beta\_3 \to \beta\_2) =
  \mathbf{inst}(f,\, \beta\_3 \to \beta\_2)\)
  e \(\mathcal{G}(\mathbf{true}, \beta\_3) = \beta\_3 =
  \mathbf{Bool}\), a restrição torna-se:

  \[ \exists \beta_1\, \beta_2\, \beta_3.\; \bigl(\mathbf{def}\; f :
      \beta_1\; \mathbf{in}\; \mathbf{inst}(f,\, \beta_3 \to \beta_2) \wedge
  \beta_3 = \mathbf{Bool}\bigr) \;\wedge\; (\alpha_0 = \beta_1 \to \beta_2) \]

  A resolução força \(\beta_1 = \beta_3 \to \beta_2\), combinada com
  \(\beta\_3 = \mathbf{Bool}\), produz
  \(\beta_1 = \mathbf{Bool} \to \beta_2\), e o tipo final é \(\alpha_0
  = (\mathbf{Bool} \to \beta_2) \to \beta_2\). Sem mais restrições
  sobre \(\beta_2\), o tipo é polimórfico: \(\forall
  \beta.; (\mathbf{Bool} \to \beta) \to \beta\).

  \section{O Resolvedor de Restrições}

  O resolvedor percorre a restrição em profundidade, mantendo duas
  estruturas de estado: a \textbf{substituição acumulada} \(S\) e o
  \textbf{contexto de tipos} \(\Gamma\). A cada etapa, a substituição é
  refinada por unificação, e o contexto é estendido por definições e ligações.

  \subsection{Unificação}

  O núcleo do resolvedor é o algoritmo de \textbf{unificação de
  Robinson}. Dado dois tipos \(T_1\) e \(T_2\), a unificação tenta
  encontrar uma substituição \(S\) tal que \(S(T_1) = S(T_2)\). O
  algoritmo mais geral (\(\mathrm{mgu}\)) produz o \textbf{unificador
  mais geral} (most general unifier):

  \[
    \begin{array}{rcll} \mathrm{mgu}(\alpha, \alpha) & = & \varepsilon &
      \text{variável consigo mesma} \\ \mathrm{mgu}(\alpha, T) & = &
      [\alpha \mapsto
      T] & \text{se } \alpha \notin \mathrm{fv}(T) \\ \mathrm{mgu}(T,
      \alpha) & = &
      [\alpha \mapsto T] & \text{se } \alpha \notin \mathrm{fv}(T) \\
      \mathrm{mgu}(T_1 \to T_2,\; T_1' \to T_2') & = & S_2 \circ S_1 &
      \text{onde }
      S_1 = \mathrm{mgu}(T_1, T_1'),\; S_2 = \mathrm{mgu}(S_1(T_2),
      S_1(T_2')) \\
      \mathrm{mgu}(T, T) & = & \varepsilon & \text{tipos concretos iguais} \\
      \mathrm{mgu}(T_1, T_2) & = & \text{erro} & \text{caso contrário}
  \end{array} \]

  A condição \(\alpha \notin \mathrm{fv}(T)\) é o \textbf{occurs
  check}: ela evita a criação de tipos recursivos infinitos como
  \(\alpha = \alpha \to \alpha\). Se o occurs check falhar, a restrição
  não tem solução.

  \subsection{O Algoritmo de Resolução}

  O algoritmo \(\mathrm{solve}(C)\) percorre \(C\) recursivamente:

  \[
    \begin{array}{rcl} \mathrm{solve}(\top) & = & \text{sucesso} \\
      \mathrm{solve}(T_1 = T_2) & = & \text{unifica } T_1 \text{ com } T_2,\;
      \text{estende } S \\ \mathrm{solve}(C_1 \wedge C_2) & = &
      \mathrm{solve}(C_1);\; \mathrm{solve}(C_2) \\ \mathrm{solve}(\exists
      \alpha.\; C) & = & \text{gera } \alpha \text{ é uma nova variável};\;
      \mathrm{solve}(C\,[\alpha]) \\ \mathrm{solve}(\mathbf{def}\; x : T\;
      \mathbf{in}\; C) & = & \text{estende } \Gamma \text{ com } x : T;\;
      \mathrm{solve}(C);\; \text{remove } x \text{ de } \Gamma \\
      \mathrm{solve}(\mathbf{let}\; x : T\; [C_1]\; \mathbf{in}\; C_2) & = &
      \mathrm{solve}(C_1);\; \sigma \leftarrow \mathrm{gen}(\Gamma, S(T));\;
      \text{estende } \Gamma \text{ com } x : \sigma;\; \mathrm{solve}(C_2) \\
      \mathrm{solve}(\mathbf{inst}(x,\, T)) & = & \sigma \leftarrow
      \Gamma(x);\; T'
      \leftarrow \mathrm{instancia}(\sigma);\; \mathrm{solve}(T = T') \\
      \mathrm{solve}(\mathbf{instScheme}(\sigma,\, T)) & = & \text{verifica
      compatibilidade de } S(T) \text{ com } \sigma
  \end{array} \]

  O tratamento de \(\mathbf{let}\) é o coração da \textbf{polimorfismo
  de let}: após resolver \(C_1\), o tipo \(T\) da variável \(x\) já foi
  parcialmente determinado pela substituição acumulada \(S\). A
  generalização \(\mathrm{gen}(\Gamma, S(T))\) quantifica
  universalmente as variáveis livres em \(S(T)\) que não ocorrem no
  contexto \(\Gamma\) (i.e., que não são ``comprometidas'' por
  dependências externas). O resultado é um esquema polimórfico que
  permite a \(x\) ser usado com diferentes instâncias em \(C_2\).

  A subtleza da regra para \(\mathbf{inst}\) está na
  \textbf{instanciação}: o esquema \(\sigma = \forall \vec{\alpha}.\;
  T'\) de \(x\) é instanciado substituindo \(\vec{\alpha}\) por
  variáveis frescas \(\vec{\beta}\) antes da unificação com \(T\). Isso
  permite que usos distintos de \(x\) em diferentes partes do programa
  recebam instâncias diferentes do esquema.

  \section{Correção do Algoritmo de Inferência}

  \subsection{Enunciados}

  Sejam \(e\) uma expressão fechada e \(\alpha_0\) uma variável de tipo
  fresca. Defina \(C_e = \exists \alpha_0.\; \mathcal{G}(e, \alpha_0)\).

  \textbf{Teorema (Correção).} Se o resolvedor aceita \(C_e\) com
  substituição final \(S\), então \(\vdash e : S(\alpha_0)\).

  Em outras palavras: o tipo produzido pela inferência é um tipo válido
  para \(e\) segundo as regras de tipagem de Hindley-Milner.

  \textbf{Teorema (Completude).} Se \(\Gamma \vdash e : T\) para algum
  \(\Gamma\) e \(T\), então \(C_e\) é satisfazível e o resolvedor
  aceita \(C_e\), produzindo um tipo \(S(\alpha_0)\) tal que \(T\) é
  instância de \(S(\alpha_0)\).

  Em outras palavras: se \(e\) é tipável, o algoritmo encontra um tipo,
  e esse tipo é o \textbf{tipo principal}, o mais geral possível.

  \textbf{Teorema (Tipo Principal).} Para qualquer expressão fechada
  tipável \(e\), existe um esquema de tipo principal \(\sigma\) tal que:

  \begin{enumerate}
    \item
      \(\vdash e : T\) para todo tipo \(T\) que é instância de \(\sigma\).
    \item
      Para todo \(T'\) tal que \(\vdash e : T'\), \(T'\) é instância de
      \(\sigma\).
  \end{enumerate}

  O tipo produzido pelo algoritmo, após generalizar \(S(\alpha_0)\),
  é exatamente esse tipo principal.

  \subsection{Esboços de Demonstração}

  \subsubsection{Correção}

  A demonstração é por indução sobre a estrutura da expressão \(e\),
  mostrando que se \(S \models \mathcal{G}(e, \alpha)\), então \(\vdash
  e : S(\alpha)\).

  \textbf{Caso (G-Var), \(e = x\):} A restrição é \(\mathbf{inst}(x,
  \alpha)\). O resolvedor instancia o esquema \(\sigma = \Gamma(x)\)
  com variáveis frescas e unifica com \(\alpha\). Logo \(S(\alpha)\) é
  uma instância de \(\sigma\), e a regra (T-Var) do sistema de tipos dá
  \(\Gamma \vdash x : S(\alpha)\).

  \textbf{Caso (G-Lit), \(e = \ell\):} A restrição
  \(\alpha = \mathrm{tylit}(\ell)\) é
  satisfeita se e somente se \(S(\alpha) =
  \mathrm{tylit}(\ell)\). As regras (T-Int) e (T-Bool)
  fornecem imediatamente \(\vdash \ell : \mathrm{tylit}(\ell)\).

  \textbf{Caso (G-App), \(e = e_1\; e_2\):} A restrição é
  \(\exists \beta.\;
    \mathcal{G}(e\_1,\, \beta \to \alpha) \wedge
  \mathcal{G}(e\_2,\, \beta)\). Por hipótese
  de indução: \(\vdash e\_1 : S(\beta \to \alpha) =
  S(\beta) \to S(\alpha)\) e \(\vdash e_2 : S(\beta)\). A
  regra (T-App) conclui \(\vdash e\_1\; e\_2
  : S(\alpha)\).

  \textbf{Caso (G-Lam), \(e = \lambda x.\; t\):} A restrição introduz
  \(\beta\_1, \beta\_2\) frescos, define \(x
  : \beta_1\) no contexto, e exige \(\mathcal{G}(t, \beta_2) \wedge
  (\alpha = \beta_1 \to \beta_2)\). Após a resolução, \(S(\alpha) =
  S(\beta_1) \to S(\beta_2)\). Por hipótese de indução (com \(x :
  S(\beta_1)\) no contexto): \(x : S(\beta\_1) \vdash t :
  S(\beta\_2)\). A regra (T-Abs) conclui
  \(\vdash \lambda x.\; t : S(\beta\_1) \to
  S(\beta\_2)\).

  \textbf{Caso (G-Let), \(e = \mathbf{let}\; x = e_1\; \mathbf{in}\;
  e_2\):} A restrição é \(\mathbf{let}\; x :
    \beta\; {[}\mathcal{G}(e\_1, \beta){]}\;
  \mathbf{in}\; \mathcal{G}(e\_2, \alpha)\).
  Após resolver \(\mathcal{G}(e\_1, \beta)\),
  o resolvedor generaliza \(S(\beta)\) em \(\sigma =
  \mathrm{gen}(\Gamma, S(\beta))\). Por hipótese de
  indução: \(\vdash e\_1 : S(\beta)\). A
  generalização produz o esquema mais geral; por hipótese de indução
  aplicada a \(e_2\) com \(x : \sigma\) no contexto: \(x :
  \sigma \vdash e\_2 : S(\alpha)\). A regra (T-Let) conclui.

  \subsubsection{Completude}

  A demonstração vai na direção oposta: dado que \(\Gamma \vdash e :
  T\), mostramos que o algoritmo encontra uma solução.

  A chave é a \textbf{propriedade de principalidade da unificação}: o
  algoritmo de Robinson sempre produz o unificador mais geral quando
  uma solução existe. Isso implica que, se \(T\) é um tipo válido para
  \(e\), então o tipo inferido \(S(\alpha_0)\) é ``mais geral que
  \(T\)'': existe uma substituição \(R\) tal que \(R(S(\alpha_0)) = T\).

  Por indução sobre a derivação de \(\Gamma \vdash e : T\):

  \begin{itemize}
    \item
      Para cada regra de tipagem, a regra de geração correspondente
      cria exatamente as equações que a regra de tipagem exige. Como o
      unificador de Robinson é completo (encontra solução se ela
      existir), as equações geradas para os subtermos são satisfazíveis.
    \item
      O tratamento do \(\mathbf{let}\) usa a generalização, que produz
      exatamente o esquema mais geral: generaliza todas as variáveis
      ``livres'' que não dependem do contexto externo, garantindo que
      qualquer instância usada no corpo seja derivável.
  \end{itemize}

  A demonstração completa requer lemas auxiliares sobre a preservação
  da satisfatibilidade pela composição de substituições e pela extensão
  de contextos, que seguem de propriedades padrão da unificação.

  \section{Eliminação de Tipos}

  A inferência de tipos produz, além do tipo principal de uma expressão,
  uma árvore sintática \emph{elaborada} cujos nós carregam anotações de
  tipo. Essas anotações são indispensáveis durante a análise estática, mas
  são desnecessárias em tempo de execução: em sistemas do estilo
  Hindley-Milner, os tipos não influenciam a semântica operacional. A
  \textbf{eliminação de tipos} (\emph{type erasure}) é a operação que
  remove todas as anotações, traduzindo a expressão elaborada de volta para
  o $\lambda$-cálculo não tipado. Isso permite reutilizar integralmente o
  \emph{pipeline} de compilação do $\lambda$-cálculo para código de máquina.

  \subsection{A Função de Eliminação}

  A eliminação de tipos é definida como uma função \(\lfloor \cdot \rfloor :
  \mathit{TyExp} \to \mathit{Term}\) que descarta as anotações de tipo e
  preserva a estrutura computacional do termo.

  \[
    \begin{array}{rcl}
      \lfloor x : T \rfloor & = & x \\[4pt]
      \lfloor \ell : T \rfloor & = & \ell \\[4pt]
      \lfloor (te_1\; te_2) : T \rfloor & = & \lfloor te_1 \rfloor\;
        \lfloor te_2 \rfloor \\[4pt]
      \lfloor (\lambda x {:} T_1.\; te) : T \rfloor & = & \lambda x.\;
        \lfloor te \rfloor \\[4pt]
      \lfloor \mathbf{let}\; x = te_1\; \mathbf{in}\; te_2 : T \rfloor & = &
        (\lambda x.\; \lfloor te_2 \rfloor)\; \lfloor te_1 \rfloor
    \end{array}
  \]

  Os quatro primeiros casos são diretos: variáveis, literais, aplicações e
  abstrações perdem apenas as anotações de tipo, mantendo a mesma estrutura
  de árvore. O caso do \haskell{let} merece atenção especial.

  \textbf{Eliminação do \texttt{let}.} A construção \(\mathbf{let}\; x =
  te_1\; \mathbf{in}\; te_2\) não existe no $\lambda$-cálculo não tipado. A
  eliminação a traduz para um \textbf{beta-redex}: \((\lambda x.\; \lfloor
  te_2 \rfloor)\; \lfloor te_1 \rfloor\). Isso é semanticamente correto
  porque, pela regra beta, \((\lambda x.\; t_2)\; t_1 \leadsto
  [x \mapsto t_1]\, t_2\), que corresponde exatamente à semântica do
  \haskell{let}. A diferença entre \haskell{let} e lambda é, portanto,
  unicamente estática: o sistema de tipos trata \(\mathbf{let}\) com
  generalização (polimorfismo de let), mas a semântica de ambas as formas é
  idêntica após a eliminação.

  \subsection{Propriedades da Eliminação}

  A principal propriedade da eliminação de tipos é a
  \textbf{preservação semântica}: avaliar o termo tipado e apagar as
  anotações do resultado produz o mesmo valor que apagar as anotações
  primeiro e avaliar o termo não tipado.

  Para enunciar este teorema precisamos relacionar os valores tipados com
  os valores não tipados. Seja \(\lfloor v \rfloor\) a extensão natural de
  \(\lfloor \cdot \rfloor\) a valores: \(\lfloor n \rfloor = n\),
  \(\lfloor b \rfloor = b\), e \(\lfloor \langle \lambda x.\; te,\, \rho
  \rangle \rfloor = \langle \lambda x.\; \lfloor te \rfloor,\, \lfloor
  \rho \rfloor \rangle\), onde \(\lfloor \rho \rfloor\) apaga as anotações
  em cada valor do ambiente.

  \begin{Theorem}[Preservação semântica da eliminação]
    Sejam \(te\) uma expressão elaborada fechada e \(v\) um valor tipado.
    Se \(\,te \Downarrow v\) na semântica tipada, então \(\lfloor te
    \rfloor \Downarrow \lfloor v \rfloor\) na semântica do $\lambda$-cálculo
    não tipado.
  \end{Theorem}

  A demonstração é por indução sobre a derivação de \(te \Downarrow v\). Os
  casos de variável, literal, abstração e aplicação são imediatos da
  correspondência estrutural entre as duas semânticas. O caso do
  \haskell{let} segue do fato de que a semântica tipada avalia \(te_1\)
  em \(v_1\) e então \([x \mapsto v_1]\,te_2\) em \(v\), enquanto o
  beta-redex avalia \(\lfloor te_1 \rfloor\) em \(\lfloor v_1 \rfloor\) e
  depois \([x \mapsto \lfloor v_1 \rfloor]\,\lfloor te_2 \rfloor\) em
  \(\lfloor v \rfloor\), que coincide com a hipótese de indução.

  Uma consequência imediata é que a eliminação de tipos é segura: nenhuma
  informação computacionalmente relevante é perdida. Os tipos existem
  exclusivamente para a análise estática e, uma vez que o programa é
  considerado bem tipado, podem ser completamente descartados.

  \section{Implementação em Haskell}

  A implementação da inferência de tipos para MiniML divide-se em cinco
  módulos: a representação de tipos e esquemas
  (\haskell{Type.hs}), a linguagem de restrições
  (\haskell{Constraint.hs}), o gerador de restrições
  (\haskell{ConstraintGen.hs}), o resolvedor
  (\haskell{Solver.hs}), e o módulo de elaboração
  (\haskell{ElabGen.hs}), que produz simultaneamente
  restrições e árvores tipadas.

  \subsection{Tipos e Esquemas}

  O módulo \haskell{Type.hs} representa tipos e substituições:

\begin{minted}{haskell}
data Type
  = TInt
  | TBool
  | TFun Type Type
  | TVar Name

data Scheme = Forall [Name] Type

type Subst = Map Name Type

-- Composição de substituições
(@@) :: Subst -> Subst -> Subst
s1 @@ s2 = Map.map (apply s1) s2 `Map.union` s1
\end{minted}

  A classe \haskell{Apply} define a aplicação de
  substituições a qualquer estrutura que contenha variáveis de tipo:

\begin{minted}{haskell}
class Apply a where
  ftv   :: a -> Set Name
  apply :: Subst -> a -> a

instance Apply Type where
  apply s (TVar n) = maybe (TVar n) id (Map.lookup n s)
  apply s (TFun t1 t2) = TFun (apply s t1) (apply s t2)
  apply _ t = t

instance Apply Scheme where
  apply s (Forall vs t) = Forall vs (apply s' t)
    where s' = Map.filterWithKey (\v _ -> v `notElem` vs) s
\end{minted}

  \subsection{Restrições}

  O módulo \haskell{Constraint.hs} representa a
  linguagem de restrições como um tipo de dados Haskell, correspondendo
  diretamente à gramática da seção anterior:

\begin{minted}{haskell}
data Constraint
  = CTrue
  | Type :=: Type
  | Constraint :&: Constraint
  | CExists (Type -> Constraint)
  | CDef Name Type Constraint
  | CLet Name Type Constraint Constraint
  | CInst Name Type
  | CInstScheme Scheme Type
\end{minted}

  O construtor \haskell{CExists} usa uma \textbf{função
  de ordem superior} para representar o escopo existencial: em vez de
  nomear explicitamente a variável, recebe uma função
  \haskell{Type -> Constraint} que, ao ser aplicada a
  um tipo, produz a restrição interna. Isso evita a necessidade
  de gerenciar nomes de variáveis na gramática de restrições e delega a
  criação de novas variáveis ao resolvedor.

  \subsection{O Gerador de Restrições}

  O módulo \haskell{ConstraintGen.hs} implementa a
  função \(\mathcal{G}(e, \alpha)\) descrita na seção anterior. Cada
  equação corresponde a uma regra de geração:

\begin{minted}{haskell}
generator :: Exp -> Type -> Constraint
generator (Var n) ty = CInst n ty
generator (Lit l) ty = ty :=: typeOfLit l
generator (App e1 e2) ty
  = CExists (\ty' ->
      generator e1 (TFun ty' ty) :&: generator e2 ty')
generator (Lam x mt body) ty
  = CExists (\t1 -> CExists (\t2 ->
      let cBody = generator body t2
          cBase = CDef x t1 cBody :&: (ty :=: TFun t1 t2)
      in maybe cBase (\t -> cBase :&: (t :=: t1)) mt))
generator (Let n mAnn e body) ty
  = CExists (\t1 ->
      let c1    = generator e t1
          c1Ann = maybe c1 (\sch -> c1 :&: CInstScheme sch t1) mAnn
          cBody = generator body ty
      in CLet n t1 c1Ann cBody)
\end{minted}

  O ponto de entrada cria a variável raiz \(\alpha_0\) via
  \haskell{CExists} e invoca o resolvedor:

\begin{minted}{haskell}
infer :: Exp -> Either String (Constraint, Env, Type)
infer e
  = let c = CExists (\t -> generator e t)
    in case solver c of
         Left err       -> Left err
         Right (env, t) -> Right (c, env, t)
\end{minted}

  \subsection{O Resolvedor}

  O módulo \haskell{Solver.hs} implementa
  \(\mathrm{solve}(C)\) na mônada \haskell{Solve = ExceptT String
  (State Env)}, onde \haskell{Env}
  mantém a substituição acumulada, o contexto de tipos, um contador
  para novas variáveis e um registro dos esquemas inferidos para
  ligações \haskell{let}.

\begin{minted}{haskell}
solve :: Constraint -> Solve ()
solve CTrue = pure ()
solve (t1 :=: t2) = do
    s  <- askSubst
    s' <- mgu (apply s t1) (apply s t2)
    extSubst s'
solve (c1 :&: c2) = solve c1 >> solve c2
solve (CExists c) = do
    v <- freshTyVar
    solve (c v)
solve (CDef x t c) =
    withLocalCtx x (Forall [] t) (solve c)
solve (CLet x t c1 c2) = do
    solve c1
    scheme <- generalize t
    addInferedType x scheme
    withLocalCtx x scheme (solve c2)
solve (CInst x t) = do
    t' <- lookupVar x
    solve (t :=: t')
solve (CInstScheme sch inf) = do
    s    <- askSubst
    tann <- instantiate sch
    ann  <- generalize tann
    inf1 <- generalize (apply s inf)
    unless (pretty ann == pretty inf1) $
      throwError "Annotated type is not an instance of the inferred type"
\end{minted}

  A função \haskell{generalize} implementa
  \(\mathrm{gen}(\Gamma, S(T))\):

\begin{minted}{haskell}
generalize :: Type -> Solve Scheme
generalize t = do
    env  <- askEnv
    s    <- askSubst
    modify (\c -> c{ctx = apply s (ctx c)})
    let t1   = apply s t
        vars = Set.toList (Set.difference (ftv t1) (free env))
    pure (Forall vars t1)
\end{minted}

  As variáveis quantificadas são exatamente
  \(\mathrm{fv}(S(T)) \setminus \mathrm{fv}(\Gamma)\).

  A unificação de Robinson é implementada por \haskell{mgu}:

\begin{minted}{haskell}
mgu :: Type -> Type -> Solve Subst
mgu t1@(TVar v) t
  | t1 == t = pure Map.empty
  | occurs v t = throwError "Occurs check failed"
  | otherwise = pure (Map.singleton v t)
mgu t (TVar v) = mgu (TVar v) t
mgu (TFun l r) (TFun l' r') = do
    s1 <- mgu l l'
    s2 <- mgu (apply s1 r) (apply s1 r')
    pure (s2 @@ s1)
mgu t1 t2
  | t1 == t2 = pure Map.empty
  | otherwise = throwError ("Cannot unify: " ++ show t1 ++ " != " ++ show t2)
\end{minted}

  \subsection{Árvores Tipadas e Elaboração}

  Uma limitação do pipeline acima é que o resultado da inferência é
  apenas um tipo: a árvore sintática original permanece sem anotações.
  Para produzir uma \textbf{árvore tipada} onde cada nó carrega o seu tipo.

  A ideia central é que o gerador produz, além da restrição, uma árvore
  tipada com \textbf{variáveis de tipo como marcadores de posição}.
  Após a resolução, a aplicação da substituição final substitui todos
  os marcadores pelos tipos concretos. Dessa forma,
  a árvore tipada é construída composicionalmente durante a geração,
  sem precisar aguardar os resultados do resolvedor; o resolvedor opera
  independentemente; e a substituição final conecta os dois.

  \subsubsection{A Árvore Tipada}

  O módulo \haskell{TyExp.hs} define a árvore de
  sintaxe abstrata com anotações de tipo em cada nó:

\begin{minted}{haskell}
data TyExp
  = TEVar  Name  Type
  | TELit  Lit   Type
  | TEApp  TyExp TyExp  Type
  | TELam  Name  Type   TyExp  Type
  | TELet  Name  TyExp  TyExp  Type
\end{minted}

  Como \haskell{TyExp} contém tipos em cada nó, a
  classe \haskell{Apply} pode ser instanciada para ela,
  permitindo que a substituição final seja aplicada à árvore inteira de
  uma só vez:

\begin{minted}{haskell}
instance Apply TyExp where
  apply s (TEVar n t) = TEVar n (apply s t)
  apply s (TELit l t) = TELit l (apply s t)
  apply s (TEApp e1 e2 t) = TEApp (apply s e1) (apply s e2) (apply s t)
  apply s (TELam x t1 e t) = TELam x (apply s t1) (apply s e) (apply s t)
  apply s (TELet n e1 e2 t) = TELet n (apply s e1) (apply s e2) (apply s t)
\end{minted}

  \subsubsection{O Gerador de Elaboração}

  O módulo \haskell{ElabGen.hs} reimplementa o gerador
  de restrições em uma mônada \haskell{State Int} que
  aloca variáveis frescas \textbf{durante a geração} (em vez de delegar
  ao resolvedor via \haskell{CExists}). As variáveis
  geradas recebem nomes \(a\_0, a\_1,
  \ldots\) para não conflitar com as variáveis \(v_0, v_1,
  \ldots\) do resolvedor.

\begin{minted}{haskell}
type ElabM = State Int

freshVar :: ElabM Type
freshVar = do { n <- get; modify (+1); return (TVar ("a" ++ show n)) }

elaborate :: Exp -> Type -> ElabM (Constraint, TyExp)
elaborate (Var n) ty
  = return (CInst n ty, TEVar n ty)
elaborate (Lit l) ty
  = return (ty :=: typeOfLit l, TELit l ty)
elaborate (App e1 e2) ty = do
    ty' <- freshVar
    (c1, te1) <- elaborate e1 (TFun ty' ty)
    (c2, te2) <- elaborate e2 ty'
    return (c1 :&: c2, TEApp te1 te2 ty)
elaborate (Lam x mt body) ty = do
    t1 <- freshVar
    t2 <- freshVar
    (cBody, teBody) <- elaborate body t2
    let cBase = CDef x t1 cBody :&: (ty :=: TFun t1 t2)
        cAnn  = maybe cBase (\t -> cBase :&: (t :=: t1)) mt
    return (cAnn, TELam x t1 teBody ty)
elaborate (Let n mAnn e body) ty = do
    t1 <- freshVar
    (c1, te1)       <- elaborate e t1
    let c1Ann = maybe c1 (\sch -> c1 :&: CInstScheme sch t1) mAnn
    (cBody, teBody) <- elaborate body ty
    return (CLet n t1 c1Ann cBody, TELet n te1 teBody ty)
\end{minted}

  \subsubsection{O Pipeline de Elaboração}

  A função \haskell{inferElab} integra a
  geração, resolução e aplicação da substituição:

\begin{minted}{haskell}
inferElab :: Exp -> Either String (Constraint, Env, TyExp, Type)
inferElab e
  = let rootVar      = TVar "a0"
        ((c, te), _) = runState (elaborate e rootVar) 1
    in case runSolveOnly (solve c) of
         Left err  -> Left err
         Right env ->
           let s = subst env
           in Right (c, env, apply s te, apply s rootVar)
\end{minted}

  O fluxo é:

  \begin{enumerate}
    \item
      \textbf{Geração:} \haskell{elaborate e rootVar}
      produz uma restrição \(C\) e uma árvore
      \haskell{te} onde cada tipo é uma variável
      \(a_i\) (marcador de posição). O contador começa em 1 porque
      \(a_0\) é reservado como variável raiz.
    \item
      \textbf{Resolução:} \haskell{solve c} percorre
      \(C\), realizando unificações e generalizações. O resultado é um
      ambiente \haskell{env} com a substituição final
      \(S = \mathtt{subst\; env}\).
    \item
      \textbf{Elaboração:} \haskell{apply s te} aplica
      \(S\) a todos os marcadores de posição em
      \haskell{te}, produzindo a árvore tipada final.
      \haskell{apply s rootVar} retorna o tipo principal de \(e\).
  \end{enumerate}

  A arquitetura em duas fases, geração separada de resolução, é a
  característica central da abordagem HM(X): o gerador é uma função
  pura que não precisa saber nada sobre unificação, e o resolvedor não
  precisa saber nada sobre a linguagem. A substituição final é o único
  ponto de contato entre os dois.

  \subsection{Eliminação de Tipos e Compilação para o $\lambda$-cálculo}

  Uma vez que \haskell{inferElab} produz a árvore tipada
  \haskell{te :: TyExp} com todos os tipos concretos, a compilação
  prossegue em duas etapas adicionais: a eliminação de tipos, que
  traduz \haskell{te} para um termo do $\lambda$-cálculo não tipado, e a
  compilação desse termo para a linguagem de destino, reutilizando o
  \emph{pipeline} já existente.

  \subsubsection{O Módulo de Eliminação}

  O módulo \haskell{MiniML.Backend.Lambda.Erase} implementa a função
  \(\lfloor \cdot \rfloor\) da seção anterior. Cada equação da definição
  matemática corresponde diretamente a uma equação em Haskell:

\begin{minted}{haskell}
module MiniML.Backend.Lambda.Erase (erase) where

import qualified Lambda.Frontend.Syntax.Term as L
import MiniML.Frontend.Syntax.Exp (Lit (..))
import MiniML.Frontend.Syntax.TyExp

erase :: TyExp -> L.Term
erase (TEVar x _)       = L.Var x
erase (TELit l _)       = L.Lit (eraseLit l)
erase (TEApp e1 e2 _)   = L.App (erase e1) (erase e2)
erase (TELam x _ e _)   = L.Lam x (erase e)
erase (TELet n e1 e2 _) = L.App (L.Lam n (erase e2)) (erase e1)

eraseLit :: Lit -> L.Lit
eraseLit (LInt n)  = L.LInt n
eraseLit (LBool b) = L.LBool b
\end{minted}

  As anotações de tipo de cada construtor (o último argumento) são
  simplesmente descartadas pelo padrão \haskell{_}. O caso do
  \haskell{let}, na última equação, gera o beta-redex
  \texttt{App (Lam n (erase e2)) (erase e1)}, conforme a definição
  formal da seção anterior.

  \subsubsection{O Pipeline Completo}

  Com o módulo de eliminação, o compilador de MiniML para TImp conecta
  quatro fases em sequência:

  \begin{enumerate}
    \item \textbf{Análise sintática} (\haskell{parser}): transforma a
      cadeia de entrada em uma árvore \haskell{Exp}.
    \item \textbf{Inferência e elaboração} (\haskell{inferElab}): produz
      a árvore tipada \haskell{TyExp} com todos os tipos concretos.
    \item \textbf{Eliminação} (\haskell{erase}): descarta as anotações
      de tipo, produzindo um \haskell{Term} do $\lambda$-cálculo não tipado.
    \item \textbf{Compilação para TImp} (\haskell{lambdaToTImp}): aplica
      o \emph{pipeline} de \emph{closure conversion} e promoção de
      lambdas para gerar um programa TImp.
  \end{enumerate}

  O comando \haskell{:timp} do REPL de MiniML executa esse pipeline
  completo:

\begin{minted}{haskell}
processTImp :: String -> IO ()
processTImp input =
  case parser (dropWhile (== ' ') input) of
    Left err  -> putStrLn $ "Parse error:\n" ++ err
    Right ast ->
      case inferElab ast of
        Left err -> putStrLn $ "Type error: " ++ err
        Right (_, _, te, ty) -> do
          putStrLn $ "val : " ++ pretty ty
          let lterm = erase te
              prog  = lambdaToTImp lterm
          result <- interpret prog
          case result of
            Left err -> putStrLn $ "Runtime error: " ++ err
            Right () -> return ()
\end{minted}

  Observe que o tipo inferido é exibido \emph{antes} da execução: ele
  provém da fase de inferência, que já terminou. A execução em si opera
  sobre o $\lambda$-cálculo não tipado, sem qualquer referência a tipos em
  tempo de execução, confirmando o princípio da eliminação.

  \textbf{Exemplo.} Para a entrada \haskell{let f = \x . x in f 42}:

  \begin{enumerate}
    \item A inferência produz o tipo \haskell{int} e a árvore tipada
      \(te = \mathbf{let}\; f = (\lambda x{:}\alpha.\; x{:}\alpha){:}\alpha
      \to \alpha\; \mathbf{in}\; (f\; 42){:}\mathbf{Int}\).
    \item A eliminação produz o lambda-termo
      \((\lambda f.\; f\; 42)\; (\lambda x.\; x)\).
    \item A \emph{closure conversion} promove \(\lambda f\) e \(\lambda
      x\) em funções de topo, representando cada closure como um registro
      \texttt{Closure}.
    \item O intérprete TImp avalia o programa e imprime \texttt{42}.
  \end{enumerate}

  A modularidade dessa arquitetura é a sua maior vantagem: a fase de
  inferência de tipos não precisa conhecer o formato de saída, e o
  \emph{backend} de $\lambda$-cálculo não precisa saber nada sobre sistemas
  de tipos. A eliminação de tipos é o ponto de desacoplamento que une as
  duas partes do compilador.

  \section{Conclusão}

  Este capítulo desenvolveu a inferência de tipos para MiniML por meio
  da abordagem de geração e resolução de restrições. A função
  \(\mathcal{G}(e, \alpha)\) transforma o
  problema de inferência em um sistema de equações entre tipos que
  captura exatamente as condições que qualquer atribuição de tipos
  válida deve satisfazer. O resolvedor de restrições, baseado na
  unificação de Robinson e na generalização de Milner, encontra a
  solução mais geral possível: o tipo principal.

  Os teoremas de correção e completude garantem que o algoritmo não
  produz falsos positivos nem falsos negativos: se o algoritmo aceita
  um programa, ele é verdadeiramente bem tipado; se o algoritmo
  rejeita, não existe nenhuma atribuição de tipos válida. O tipo
  principal captura toda a informação de tipo disponível para um
  programa: qualquer tipo válido é uma instância do tipo principal.

  A implementação em Haskell evidencia a correspondência sistemática
  entre a especificação matemática e o código. Cada regra de geração
  corresponde a uma equação do gerador; cada linha do resolvedor
  implementa um caso da definição semântica de satisfazibilidade; e a
  classe \haskell{Apply} é a realização computacional
  da substituição matemática. A separação em módulos (tipos,
  restrições, gerador, resolvedor, elaborador) reflete a
  modularidade conceitual do algoritmo.

  A extensão para produzir árvores tipadas, seguindo a elaboração em
  estilo aplicativo, mostra que a inferência de tipos não precisa se
  limitar a produzir um tipo: ela pode ser vista como uma transformação
  que anota o programa com toda a informação de tipo descoberta durante
  a inferência, servindo de base para fases posteriores do compilador,
  como geração de código especializado ou análises de otimização
  sensíveis a tipos.
